% Môn: Vật lý
% Lớp: 12
% Chương: Chương I. Vật lí nhiệt
% Bài: L12C1 Bài 1. Cấu trúc của chất. Sự chuyển thể
% Loại: Ôn tập
% File riêng ôn tập bài 1

% =========================================================
% PHẦN I. CÂU TRẮC NGHIỆM NHIỀU LỰA CHỌN
% =========================================================

\begin{ex}
Tìm phát biểu \textbf{sai}.
\choice
{Nội năng là một dạng năng lượng nên có thể chuyển hóa thành các dạng năng lượng khác.}
{Nội năng của một vật phụ thuộc vào nhiệt độ và thể tích của vật.}
{\True Nội năng chính là nhiệt lượng của vật.}
{Nội năng của vật có thể tăng hoặc giảm.}
\loigiai{
Nội năng của vật không phải là nhiệt lượng của vật.
}
\end{ex}

\begin{ex}
Cách nào sau đây \textbf{không làm thay đổi nội năng của vật}?
\choice
{Cọ xát vật lên mặt bàn.}
{Đốt nóng vật.}
{Làm lạnh vật.}
{\True Đưa vật lên cao.}
\loigiai{
Đưa vật lên cao làm thay đổi thế năng của vật nhưng không làm thay đổi nội năng của vật.
}
\end{ex}

\begin{ex}
Trường hợp làm biến đổi nội năng \textbf{không do thực hiện công} là
\choice
{\True Đun nóng nước bằng bếp.}
{Một viên bi bằng thép rơi xuống đất mềm.}
{Nén khí trong xilanh.}
{Cọ xát hai vật vào nhau.}
\loigiai{
Đun nóng nước bằng bếp là quá trình truyền nhiệt, không phải thực hiện công.
}
\end{ex}

\begin{ex}
Trường hợp nội năng của vật bị biến đổi \textbf{không phải do truyền nhiệt} là
\choice
{Chậu nước để ngoài nắng một lúc nóng lên.}
{Gió mùa đông bắc tràn về làm cho không khí lạnh đi.}
{\True Khi trời lạnh, ta xoa hai bàn tay vào nhau cho ấm lên.}
{Cho cơm nóng vào bát thì thấy bát cũng nóng.}
\loigiai{
Khi xoa hai bàn tay vào nhau, nội năng tăng do thực hiện công.
}
\end{ex}

\begin{ex}
Hình bên là các dụng cụ đo nhiệt dung riêng của nước. Dụng cụ số (4) là
\choice
{Nhiệt lượng kế.}
{Nhiệt kế điện tử.}
{Cân điện tử.}
{\True Biến thế nguồn.}
\loigiai{
Dụng cụ số (4) là biến thế nguồn.
}
\end{ex}

\begin{ex}
Sự truyền nhiệt là
\choice
{Sự chuyển hóa năng lượng từ dạng này sang dạng khác.}
{\True Sự truyền trực tiếp nội năng từ vật này sang vật khác.}
{Sự chuyển hóa năng lượng từ nội năng sang dạng khác.}
{Sự truyền trực tiếp nội năng và chuyển hóa năng lượng từ dạng này sang dạng khác.}
\loigiai{
Sự truyền nhiệt là sự truyền năng lượng từ vật này sang vật khác do có sự chênh lệch nhiệt độ.
}
\end{ex}

\begin{ex}
Nhiệt lượng trao đổi trong quá trình truyền nhiệt không phụ thuộc vào
\choice
{\True Thời gian truyền nhiệt.}
{Độ biến thiên nhiệt độ.}
{Khối lượng của chất.}
{Nhiệt dung riêng của chất.}
\loigiai{
Nhiệt lượng được xác định bởi
\[
Q=mc\Delta t.
\]
Do đó không phụ thuộc trực tiếp vào thời gian truyền nhiệt.
}
\end{ex}

\begin{ex}
Đơn vị của nhiệt dung riêng của vật là
\choice
{\True $\mathrm{J/(kg\cdot K)}$.}
{$\mathrm{kg/(J\cdot K)}$.}
{$\mathrm{J\cdot kg\cdot K}$.}
{$\mathrm{kg\cdot K/J}$.}
\end{ex}

\begin{ex}
Nhiệt dung riêng của rượu là $2500\,\mathrm{J/(kg\cdot K)}$. Điều đó có nghĩa là gì?
\choice
{Để nâng $1\,\mathrm{kg}$ rượu lên nhiệt độ bay hơi phải cung cấp cho nó một nhiệt lượng $2500\,\mathrm{J}$.}
{$1\,\mathrm{kg}$ rượu bị đông đặc thì giải phóng nhiệt lượng là $2500\,\mathrm{J}$.}
{\True Để nâng $1\,\mathrm{kg}$ rượu tăng thêm $1\,\mathrm{K}$ cần cung cấp cho nó nhiệt lượng $2500\,\mathrm{J}$.}
{Nhiệt lượng có trong $1\,\mathrm{kg}$ rượu ở nhiệt độ bình thường là $2500\,\mathrm{J}$.}
\loigiai{
Theo định nghĩa nhiệt dung riêng:
\[
c=\frac{Q}{m\Delta T}.
\]
Với $m=1\,\mathrm{kg}$ và $\Delta T=1\,\mathrm{K}$ thì $Q=2500\,\mathrm{J}$.
}
\end{ex}

\begin{ex}
Gọi $t$ là nhiệt độ lúc sau, $t_0$ là nhiệt độ lúc đầu của vật. Công thức nào là công thức tính nhiệt lượng mà vật thu vào?
\choice
{$Q=mc(t_0-t)$.}
{$Q=mc(t_0+t)$.}
{\True $Q=mc(t-t_0)$.}
{$Q=mc$.}
\loigiai{
Khi vật thu nhiệt và nhiệt độ tăng:
\[
Q=mc(t-t_0).
\]
}
\end{ex}

\begin{ex}
Có 4 bình $A$, $B$, $C$, $D$ đều đựng nước ở cùng một nhiệt độ với thể tích tương ứng là $1$ lít, $2$ lít, $3$ lít và $4$ lít. Sau khi dùng các đèn cồn giống hệt nhau để đun các bình trong cùng một khoảng thời gian. Hỏi bình nào có nhiệt độ thấp nhất?
\choice
{Bình $A$.}
{Bình $B$.}
{Bình $C$.}
{\True Bình $D$.}
\loigiai{
Các bình nhận nhiệt lượng như nhau. Vì
\[
Q=mc\Delta T,
\]
bình có khối lượng nước lớn nhất có độ tăng nhiệt độ nhỏ nhất. Bình $D$ chứa nhiều nước nhất nên có nhiệt độ thấp nhất.
}
\end{ex}

\begin{ex}
Nhiệt dung riêng của đồng lớn hơn chì. Vì vậy để tăng nhiệt độ của $1\,\mathrm{kg}$ đồng và $3\,\mathrm{kg}$ chì thêm $15^\circ\mathrm{C}$ thì
\choice
{Khối chì cần nhiều nhiệt lượng hơn khối đồng.}
{\True Khối đồng cần nhiều nhiệt lượng hơn khối chì.}
{Hai khối đều cần nhiệt lượng như nhau.}
{Không khẳng định được.}
\loigiai{
Ta có
\[
Q=mc\Delta T.
\]
Với cùng $\Delta T$, khối đồng cần nhiều nhiệt lượng hơn vì $m_{\rm đồng}c_{\rm đồng}>m_{\rm chì}c_{\rm chì}$.
}
\end{ex}

\begin{ex}
Nội dung nguyên lí I nhiệt động lực học là
\choice
{\True Độ biến thiên nội năng của vật bằng tổng công và nhiệt lượng mà vật nhận được.}
{Độ biến thiên nội năng của vật bằng tổng nhiệt lượng mà vật nhận được.}
{Độ biến thiên nội năng của vật bằng tổng công mà vật nhận được.}
{Độ biến thiên nội năng của vật bằng hiệu số công và nhiệt lượng mà vật nhận được.}
\loigiai{
Theo nguyên lí I nhiệt động lực học:
\[
\Delta U=A+Q.
\]
}
\end{ex}

\begin{ex}
Biểu thức diễn tả đúng quá trình chất khí vừa nhận nhiệt vừa nhận công là
\choice
{$\Delta U=A+Q;\ Q>0;\ A<0$.}
{$\Delta U=Q;\ Q>0$.}
{$\Delta U=Q+A;\ Q<0;\ A>0$.}
{\True $\Delta U=Q+A;\ Q>0;\ A>0$.}
\loigiai{
Khí nhận nhiệt nên $Q>0$, nhận công nên $A>0$. Do đó
\[
\Delta U=Q+A.
\]
}
\end{ex}

\begin{ex}
Nội năng của vật phụ thuộc vào
\choice
{\True Nhiệt độ và thể tích của vật.}
{Khối lượng và nhiệt độ của vật.}
{Khối lượng và thể tích của vật.}
{Khối lượng của vật.}
\loigiai{
Nội năng của một vật phụ thuộc vào nhiệt độ và thể tích của vật.
}
\end{ex}

\begin{ex}
Hiện tượng quả bóng bàn bị móp (nhưng chưa bị thủng) khi thả vào cốc nước nóng sẽ phồng trở lại là do
\choice
{\True Nội năng của chất khí tăng lên.}
{Nội năng của chất khí giảm xuống.}
{Nội năng của chất khí không thay đổi.}
{Nội năng của chất khí bị mất đi.}
\loigiai{
Nước nóng truyền nhiệt cho chất khí bên trong bóng, làm nội năng của khí tăng, khí nở ra và làm bóng phồng trở lại.
}
\end{ex}

\begin{ex}
Cung cấp cho vật một công là $200\,\mathrm{J}$ nhưng nhiệt lượng bị thất thoát ra môi trường bên ngoài là $120\,\mathrm{J}$. Nội năng của vật
\choice
{\True Tăng $80\,\mathrm{J}$.}
{Giảm $80\,\mathrm{J}$.}
{Không thay đổi.}
{Giảm $320\,\mathrm{J}$.}
\loigiai{
Vật nhận công $A=200\,\mathrm{J}$ và tỏa nhiệt nên $Q=-120\,\mathrm{J}$. Do đó
\[
\Delta U=A+Q=200-120=80\,\mathrm{J}.
\]
}
\end{ex}

\begin{ex}
Hệ thức $\Delta U=A+Q$ khi $Q<0$ và $A>0$ mô tả quá trình
\choice
{Hệ truyền nhiệt và sinh công.}
{Hệ nhận nhiệt và sinh công.}
{\True Hệ truyền nhiệt và nhận công.}
{Hệ nhận nhiệt và nhận công.}
\loigiai{
Vì $Q<0$ nên hệ truyền nhiệt ra môi trường. Vì $A>0$ nên hệ nhận công.
}
\end{ex}

% =========================================================
% PHẦN II. CÂU TRẮC NGHIỆM ĐÚNG / SAI
% =========================================================

\begin{ex}
Xét khối khí như trong hình. Dùng tay ấn mạnh và nhanh pít-tông, vừa nung nóng khí.
\choiceTF
{\True Khí nhận công nên $A>0$.}
{Nhiệt lượng $Q<0$ vì khi bị nung nóng khí nhận nhiệt.}
{\True Nội năng của khí tăng, $\Delta U>0$.}
{Biểu thức liên hệ giữa độ biến thiên nội năng, công và nhiệt lượng là $\Delta U=A-Q$.}
\loigiai{
Khí bị nén nên khí nhận công, do đó $A>0$. Khí được nung nóng nên nhận nhiệt, do đó $Q>0$. Vì cả $A>0$ và $Q>0$ nên
\[
\Delta U=A+Q>0.
\]
}
\end{ex}

\begin{ex}
Trong quá trình nóng chảy của vật rắn, xét các phát biểu sau.
\choiceTF
{Nhiệt được truyền vào vật rắn để làm tăng nhiệt độ của nó.}
{Động năng trung bình của các phân tử trong vật rắn tăng lên.}
{\True Nội năng của vật rắn thay đổi.}
{Tại nhiệt độ nóng chảy, nội năng không thay đổi.}
\loigiai{
Trong quá trình nóng chảy, nhiệt độ của chất không đổi nhưng nội năng vẫn thay đổi do thế năng tương tác giữa các phân tử thay đổi.
}
\end{ex}

\begin{ex}
Xét một khối khí trong bình kín bị nung nóng.
\choiceTF
{Khí truyền nhiệt ra môi trường xung quanh.}
{Công $A\ne0$ vì thể tích khí thay đổi.}
{\True Nội năng của khí tăng.}
{\True Hệ thức phù hợp với quá trình là $\Delta U=Q$, với $Q>0$.}
\loigiai{
Bình kín có thể tích không đổi nên khí không thực hiện công. Khí được nung nóng nên nhận nhiệt $Q>0$. Do đó
\[
\Delta U=Q>0.
\]
}
\end{ex}

\begin{ex}
Một lượng nước và một lượng rượu có thể tích bằng nhau được cung cấp các nhiệt lượng tương ứng là $Q_1$ và $Q_2$. Biết khối lượng riêng của nước là $1000\,\mathrm{kg/m^3}$ và của rượu là $800\,\mathrm{kg/m^3}$, nhiệt dung riêng của nước là $4200\,\mathrm{J/(kg\cdot K)}$ và của rượu là $2500\,\mathrm{J/(kg\cdot K)}$.
\choiceTF
{Nhiệt lượng để làm tăng nhiệt độ của $1\,\mathrm{kg}$ nước lên $1\,\mathrm{K}$ là $2500\,\mathrm{J}$.}
{Nhiệt lượng để làm tăng nhiệt độ của $1\,\mathrm{kg}$ rượu lên $1\,\mathrm{K}$ là $4200\,\mathrm{J}$.}
{\True Có thể dùng công thức $Q=mc(T_2-T_1)$ để tính nhiệt lượng cung cấp cho nước và rượu.}
{\True Nếu độ tăng nhiệt độ của nước và rượu bằng nhau thì $Q_1=2{,}1Q_2$.}
\loigiai{
Ta có
\[
Q=mc\Delta T.
\]
Vì hai chất có cùng thể tích nên
\[
\frac{m_{\rm nước}}{m_{\rm rượu}}=\frac{1000}{800}=1{,}25.
\]
Do đó, khi độ tăng nhiệt độ bằng nhau,
\[
\frac{Q_1}{Q_2}=\frac{1000\cdot4200}{800\cdot2500}=2{,}1.
\]
}
\end{ex}

% =========================================================
% PHẦN III. CÂU TRẮC NGHIỆM TRẢ LỜI NGẮN
% =========================================================

\begin{ex}
Một lượng khí nhận nhiệt lượng $250\,\mathrm{kJ}$ do được đun nóng; đồng thời nhận công $500\,\mathrm{kJ}$ do bị nén. Độ tăng nội năng của lượng khí là bao nhiêu kJ?
\shortans{750}
\loigiai{
Khí nhận nhiệt và nhận công nên
\[
\Delta U=Q+A=250+500=750\,\mathrm{kJ}.
\]
}
\end{ex}

\begin{ex}
Người ta thực hiện công $200\,\mathrm{J}$ để nén khí trong một xilanh. Biết khí truyền ra môi trường xung quanh nhiệt lượng $40\,\mathrm{J}$. Độ biến thiên nội năng của khí là bao nhiêu J?
\shortans{160}
\loigiai{
Khí nhận công $A=200\,\mathrm{J}$ và truyền nhiệt ra môi trường nên $Q=-40\,\mathrm{J}$. Do đó
\[
\Delta U=A+Q=200-40=160\,\mathrm{J}.
\]
}
\end{ex}

\begin{ex}
Một quả bóng khối lượng $200\,\mathrm{g}$ rơi từ độ cao $15\,\mathrm{m}$ xuống sàn và nảy lên được $10\,\mathrm{m}$. Độ biến thiên nội năng của quả bóng bằng bao nhiêu J? Lấy $g=10\,\mathrm{m/s^2}$.
\shortans{10}
\loigiai{
Phần cơ năng bị mất chuyển hóa thành nội năng:
\[
\Delta U=mg(h_1-h_2)
=0{,}2\cdot10(15-10)=10\,\mathrm{J}.
\]
}
\end{ex}

\begin{ex}
Một viên đạn đại bác có khối lượng $10\,\mathrm{kg}$ khi tới đích có vận tốc $54\,\mathrm{km/h}$. Nếu toàn bộ động năng của nó biến thành nội năng thì nhiệt lượng tỏa ra lúc va chạm vào khoảng bao nhiêu J?
\shortans{1125}
\loigiai{
Đổi đơn vị:
\[
v=54\,\mathrm{km/h}=15\,\mathrm{m/s}.
\]
Động năng của viên đạn:
\[
W_{\mathrm{đ}}=\frac12mv^2
=\frac12\cdot10\cdot15^2
=1125\,\mathrm{J}.
\]
Vậy nhiệt lượng tỏa ra là $1125\,\mathrm{J}$.
}
\end{ex}

\begin{ex}
Để xác định nhiệt độ của một lò nung, người ta đưa vào trong lò một miếng sắt có khối lượng $50\,\mathrm{g}$. Khi miếng sắt có nhiệt độ bằng nhiệt độ của lò, người ta lấy ra và thả vào một nhiệt lượng kế chứa $900\,\mathrm{g}$ nước ở nhiệt độ $17^\circ\mathrm{C}$. Khi đó nhiệt độ của nước tăng lên đến $23^\circ\mathrm{C}$. Biết nhiệt dung riêng của sắt là $460\,\mathrm{J/(kg\cdot K)}$, của nước là $4200\,\mathrm{J/(kg\cdot K)}$. Nhiệt độ của lò xấp xỉ bằng bao nhiêu K?
\shortans{1282}
\loigiai{
Nhiệt lượng nước nhận được:
\[
Q_{\rm nước}=0{,}9\cdot4200\cdot(23-17)=22680\,\mathrm{J}.
\]
Bỏ qua nhiệt lượng kế, nhiệt lượng sắt tỏa ra bằng nhiệt lượng nước thu vào:
\[
0{,}05\cdot460(T-23)=22680.
\]
Suy ra
\[
T\approx1009^\circ\mathrm{C}.
\]
Đổi sang Kelvin:
\[
T\approx1009+273=1282\,\mathrm{K}.
\]
}
\end{ex}

\begin{ex}
Tính nhiệt lượng cần thiết theo đơn vị kJ để đun $5\,\mathrm{kg}$ nước từ $15^\circ\mathrm{C}$ đến $100^\circ\mathrm{C}$ trong một thùng bằng sắt có khối lượng $1{,}5\,\mathrm{kg}$. Biết nhiệt dung riêng của nước là $4200\,\mathrm{J/(kg\cdot K)}$, của sắt là $460\,\mathrm{J/(kg\cdot K)}$.
\shortans{1843,65}
\loigiai{
Độ tăng nhiệt độ:
\[
\Delta T=100-15=85^\circ\mathrm{C}.
\]
Nhiệt lượng nước nhận:
\[
Q_1=5\cdot4200\cdot85=1785000\,\mathrm{J}.
\]
Nhiệt lượng thùng sắt nhận:
\[
Q_2=1{,}5\cdot460\cdot85=58650\,\mathrm{J}.
\]
Tổng nhiệt lượng:
\[
Q=Q_1+Q_2=1843650\,\mathrm{J}=1843{,}65\,\mathrm{kJ}.
\]
}
\end{ex}
